\section{Introduction}
\label{sec:intro}

\subsection{The problem}

A modern terrain and solid-geometry engine must satisfy four requirements
at the same time:

\begin{enumerate}[label=\textbf{R\arabic*.},leftmargin=2.2em]
\item \textbf{Authoring freedom.}  Designers describe geometry with
  implicit primitives, Boolean operations, noise, and material painters,
  not with a fixed triangle list.
\item \textbf{Adaptive detail.}  Near geometry must be finely tessellated;
  the same region observed from a kilometre away must collapse to a coarse
  representation without cracks or holes.
\item \textbf{Low host traffic.}  A terrain that streams hundreds of
  chunks per second cannot afford a GPU$\rightarrow$CPU$\rightarrow$GPU
  round trip for every vertex.
\item \textbf{Real-time rendering.}  The result must feed frustum culling,
  shadow cascades, ray tracing, vegetation placement, and volumetric
  effects without a second geometry representation.
\end{enumerate}

The engine answers with a split pipeline.  The CPU evaluates the scene's
signed distance field (SDF) and maintains an adaptive octree of corner
samples.  It never stores a triangle mesh.  The GPU receives a flattened
subset of that octree for one chunk, performs the full adaptive Surface
Nets walk, solves one dual vertex per surface cell, and writes the
resulting vertices and indices directly into the renderer's persistent
packed pools.  The host sees integer counters, not positions.

\begin{figure}[htbp]
\centering
\resizebox{\linewidth}{!}{%
\begin{tikzpicture}[x=1cm,y=1cm,font=\small,every node/.append style={text width=3.55cm}]
  \node[io]  (author)  at (0.0,3.1) {SDF scene\\primitives, CSG, brushes};
  \node[cpu] (apply)   at (4.2,3.1) {CPU shape apply\\corner sampling};
  \node[cpu] (tree)    at (8.4,3.1) {Adaptive octree\\+ SDF[8]};
  \node[cpu] (flat)    at (12.6,3.1) {Flatten chunk\\+ cross-edge capacity};
  \node[gpu] (mesh)    at (12.6,1.2) {GPU geometry pass\\indices + vertices};
  \node[io]  (pools)   at (8.4,1.2) {Packed vertex/\\index pools};
  \node[gpu] (cull)    at (4.2,1.2) {GPU cull + LoD gate};
  \node[io]  (frame)   at (0.0,1.2) {Indirect draws,\\shadows, RT};
  \draw[arrow] (author)--(apply);
  \draw[arrow] (apply)--(tree);
  \draw[arrow] (tree)--(flat);
  \draw[arrow] (flat)--(mesh);
  \draw[arrow] (mesh)--(pools);
  \draw[arrow] (pools)--(cull);
  \draw[arrow] (cull)--(frame);
\end{tikzpicture}}
\caption{End-to-end pipeline.  Blue nodes execute on the CPU; orange nodes
execute on the GPU.  The only data crossing from GPU to CPU are integer
counters used for span sizing and statistics.  Chunk geometry itself never
returns to the host.}
\label{fig:pipeline}
\end{figure}

\subsection{Why not mesh on the CPU?}

The original engine did exactly that.  \code{Octree::iterateTriangles}
walked the tree, solved QEF vertices, and appended triangles to a
\code{Geometry} object; a staging copy then uploaded the vertex and index
arrays.  The CPU path is still present as
\code{OctreeTriangleTessellator}, but it now serves only the offline
Blender exporter.  Moving the same walk to a compute shader produced three
concrete benefits:

\begin{itemize}
\item \textbf{The host never touches the vertex stream.}  A chunk's
  topology is written into a device-local span.  The host reserves the
  span, publishes a draw command, and never reads a position.
\item \textbf{The adaptive walk parallelises trivially.}  One invocation
  owns one scan cell.  All per-cell state (break lists, QEF memo, recursion
  stack) is private, so the only atomics are the index append, the overflow
  counters, and the per-node vertex-writer election.
\item \textbf{The renderer and the mesher share one memory layout.}  The
  compute pass writes the renderer's exact 64-byte \code{Vertex} format
  and chunk-local index values.  There is no conversion, compaction, or
  intermediate mesh object.
\end{itemize}

The cost is that the GPU must see exactly the octree neighbourhood the CPU
walker would have traversed: coarser neighbours for edge stitching, finer
children for the descent, the world-root path, and a deterministic
assignment of vertex slots.  The bulk of this report is about making that
flattening exact and about the single-pass GPU mesher that consumes it.

\subsection{Contributions}

This report makes the following contributions, all implemented in the
current engine:

\begin{enumerate}[label=\textbf{C\arabic*.},leftmargin=2.2em]
\item A precise description of the CPU octree as a \emph{corner-sample
  store}, including its +1-shifted LoD representation, block allocator,
  quadratic-error simplification, and threaded \code{apply} recursion.
\item A root-consistent flattening algorithm that turns a chunk's octree
  neighbourhood into eight compact storage buffers and a scan-cell list,
  with host-assigned vertex slots packed into node flags.
\item A single-pass GPU topology geometry stage: the CPU cross-edge count
  sizes both spans, and one dispatch appends the chunk-local index stream
  and writes each referenced vertex payload exactly once.
\item Per-chunk cross-edge capacity accounting and overflow growth: an
  overflow is reported without out-of-bounds writes and grows only the
  affected chunk's span.
\item An exact GPU port of the adaptive edge stitching algorithm
  (per-edge break collection at LoD boundaries), including its fixed
  capacities, overflow policy, and deterministic ordering.
\item A packed-pool span protocol that binds only a chunk's slice of a
  large pool, respects \code{maxStorageBufferRange}, handles alignment
  padding, and publishes the draw entry only when the fence signals.
\item A hole-free hierarchical LoD gate used by the raster, shadow,
  bounding-box, and water streams, plus the GPU-resident effect consumers
  (vegetation instances and lava/fire anchors) that read the same chunk
  metadata.
\item A validation-oriented account of the synchronization, lifetime, and
  determinism rules that keep the path clean across asynchronous compute,
  graphics, and transfer queues.
\end{enumerate}

\subsection{Reading guide}

\Cref{sec:foundations}--\ref{sec:cpuwalk} cover the CPU side:
SDF mathematics, octree storage, the shape recursion, simplification, and
the reference walker.  \Cref{sec:gpudesign}--\ref{sec:finalize} cover the
GPU meshing pass in increasing detail.  \Cref{sec:orch}--\ref{sec:sync}
cover host orchestration and synchronization.  \Cref{sec:render}--
\ref{sec:effects} follow the mesh into rendering and effects.
\Cref{sec:perf}--\ref{sec:validation} give measured costs and correctness
invariants.  Appendices contain the exact tables and storage layouts.

\begin{keybox}[The one-sentence summary]
The CPU decides \emph{what the surface is}; the GPU decides \emph{where
its vertices are} and writes them once, directly into the buffers the
renderer already owns.
\end{keybox}
